\chapter{Cranioplasty}

\section*{Introduction \& Clinical Indications}
Cranioplasty reconstructs a skull defect after decompressive craniectomy, trauma, tumor surgery, infection, or congenital deformity. It protects the brain, restores cranial contour, and may improve headache, cerebrospinal-fluid dynamics, or neurological function in selected patients. Materials include the patient's stored bone, titanium, porous polyethylene, hydroxyapatite, or custom implants.

The timing is individualized after brain swelling, infection, wound condition, hydrocephalus, and systemic health are assessed. A cranial defect does not automatically require immediate reconstruction.

\section*{Relevant Surgical Anatomy}
The scalp has layered blood supply and must be mobilized without devascularization. The dura, dural substitutes, venous sinuses, brain, ventricles, and previous bone edges are protected. The implant must fit the defect without compressing the brain or obstructing CSF pathways. Temporal hollowing and frontotemporal contour are important functional and cosmetic landmarks.

\section*{Preoperative Workup \& Preparation}
CT defines the defect and implant plan. Assessment includes neurological status, hydrocephalus, shunt function, seizures, infection, wound healing, nutrition, anticoagulants, and imaging for collections or mass effect. Consent covers infection, wound breakdown, seizures, hemorrhage, CSF leak, implant exposure or resorption, hydrocephalus, neurological injury, cosmetic asymmetry, and revision.

\section*{Surgical Technique \& Procedure}
The scalp is reopened through the previous incision when safe. Scar tissue is dissected, the dura and brain are protected, and the bone or implant is shaped to restore the contour. Fixation uses plates, screws, or other approved systems. Hemostasis and watertight soft-tissue closure are prioritized; drains are used selectively.

\section*{Complications \& Risks}
Complications include infection, epidural or subgaleal hematoma, seizures, CSF leak, hydrocephalus, wound breakdown, implant exposure, bone resorption, sinking or displacement, poor contour, and neurological deterioration. Infection may require implant removal and delayed reconstruction.

\section*{Postoperative Care \& Recovery}
Neurological observations, wound checks, pain control, seizure management, and imaging are individualized. Fever, increasing headache, vomiting, confusion, weakness, seizure, wound drainage, or swelling requires urgent assessment. Follow-up evaluates protection, contour, neurological function, and hydrocephalus.

\section*{Outcomes \& Prognosis}
Most appropriately selected patients obtain durable cranial protection and improved contour, but outcomes depend on the original injury, infection history, hydrocephalus, material, timing, and wound condition. Reconstruction does not reverse all neurological injury from the original disease.

\section*{References}
\begin{enumerate}
  \item Congress of Neurological Surgeons. Guidelines for the Timing of Cranioplasty After Craniectomy.
  \item American Association of Neurological Surgeons. Patient Safety and Cranioplasty Resources.
  \item National Institute for Health and Care Excellence. Head Injury: Assessment and Early Management.
\end{enumerate}

\clearpage
